% !TEX root = ../thesis.tex

\label{ch7:outlook}

This thesis argues that inductive biases should be an essential component of generative models in the molecular sciences. The work presented here explores two means of introducing physics-inspired biases: first, by encoding thermodynamic information in the prior distribution of a diffusion model, and second, by combining physical and generative models in sequence-to-structure modeling pipelines. A few paths branch from these two initial directions.

The first is to extend Thermodynamic Maps to variables other than temperature. Temperature was the natural starting point because its effect on equilibrium fluctuations can be represented cleanly through the variance of a Gaussian prior, but the logic extends to other thermodynamic variables. Separately, exponentially tilted Thermodynamic Maps have been developed, exploring similar representations of pressure and chemical potential through structured changes in the prior \citep{exptm}. Still, several questions remain open. A central challenge is how to evaluate thermodynamic generalization: benchmark systems should have known or independently computable thermodynamics and tunable phase behavior. Spin glasses, simple fluids, polymers, and well-understood molecular systems may provide useful test beds. A second challenge is determining when a heuristic representation is sufficient. The energy-dependent prior used in RNAnneal, for example, is not a direct physical temperature in the same sense as the temperature variable in Thermodynamic Maps, but it nevertheless provides a useful ranking signal for RNA conformations. Understanding the relationship between formally grounded thermodynamic encodings and practically useful heuristic encodings will be important for broadening the method while preserving its connection to thermodynamics and statistical physics.

The second is to generate latent representations of molecular data. Recent latent-variable extensions such as latent thermodynamic flows and diffusive State Predictive Information Bottleneck suggest that representation learning and generative modeling can be mutually reinforcing, especially in the low-data regime that characterizes most molecular problems \citep{latf2026,john2025dspib}. A useful latent space should separate metastable states and preserve slow dynamical processes. Conversely, a generative objective can regularize the representation by forcing it to support faithful density estimation and sample generation. Integration of representation learning with thermodynamic prior distributions developed in this thesis may yield models that simultaneously identify collective variables, estimate free-energy landscapes, and generate ensembles under unseen conditions.

A third, applications-oriented direction is to use RNAnneal to explore the structural landscapes of therapeutically relevant RNAs, especially those with limited co-evolutionary information or heterogeneous structures. Because RNAnneal learns a sequence-specific score \textit{de novo} from physics-based sampling, it does not depend on co-evolutionary information or deep phylogenetic signal in the way that popular end-to-end structure prediction models do. In principle, that should make it applicable not only to natural RNAs with limited homologous sequence information, but also to synthetic or regulatory RNAs designed for therapeutic function. Two concrete examples are VEGF-associated regulatory RNA structures, where the goal is to characterize the disordered and condition-sensitive structural landscape \citep{ray2009vegfa}, and the MALAT1 triple helix, where structurally informed conformational ensembles could guide small-molecule screening against the unusual fold \citep{brown2014malat1,donlic2018malat1}. These applications would test whether the proposed methodology can deliver insights for targets whose functional importance lies in conformational heterogeneity or noncanonical structure.

The long-term goal is not only to predict RNA structure ensembles more accurately, but to translate them into mechanistic insight. The road ahead will entail expanding the structural vocabulary of RNAnneal beyond the motifs emphasized in the presented benchmark, including pseudoknots, G-quadruplexes, triplexes, and other noncanonical folds; linking predicted states to ligandability, accessibility, and function; and incorporating experimental measurements as modeling constraints. If successful, such direction could push sequence-to-structure models toward becoming an integral part of the decision-making process for therapeutic discovery.
